\section*{Chapter \ref{ch:m2:distressed-sovereign-and-bank-capital-credit} --- Distressed, Sovereign and Bank-Capital Credit}

\begin{solution}{exo:m2:distressed-sovereign-and-bank-capital-credit:1}
A face haircut of 40\%. The new bonds are worth 47.92 at 8\%, a present-value haircut
of 52.1\%.
\end{solution}

\begin{solution}{exo:m2:distressed-sovereign-and-bank-capital-credit:2}
A series-by-series clause needs a qualified majority, typically 75\%, of each bond
separately, so a holdout with a blocking stake in one series can keep that series out.
A single-limb aggregated clause counts one vote across all the series affected, 75\% of
their total principal in the ICMA model, so no single series can block, provided every
series is offered the same terms.
\end{solution}

\begin{solution}{exo:m2:distressed-sovereign-and-bank-capital-credit:3}
Its contractual trigger, usually the bank's common equity ratio falling below a set
level, which depends on the bank's losses; and the \omterm{def:m2:distressed-sovereign-and-bank-capital-credit:pon}{point of non-viability}, which depends
on an authority's decision that the bank would fail without a write-down or public
support.
\end{solution}

\begin{solution}{exo:m2:distressed-sovereign-and-bank-capital-credit:4}
From 60.9\% to 66.7\%: the clause binds holdouts at two thirds, so the band shrinks by
8.3 points to 5.8.
\end{solution}

\begin{solution}{exo:m2:distressed-sovereign-and-bank-capital-credit:5}
It meant that paying the creditors who had accepted the exchange required paying the
holdouts in full at the same time, and let the holdouts stop the payments through the
banks that processed them; accepting creditors lost the certainty of being paid, which
made future exchanges harder to sell. Newer contracts state that the clause does not
require ratable payment, and use aggregated \omterm{def:m2:distressed-sovereign-and-bank-capital-credit:cac}{collective action clauses}.
\end{solution}

\begin{solution}{exo:m2:distressed-sovereign-and-bank-capital-credit:6}
The Greek law could change only Greek-law bonds; the English-law bonds could not be
amended without their holders' votes, and a default on them would have exposed Greece
to lawsuits abroad. The holdouts were small relative to the whole, so paying them in
full was cheaper than the fight: the free-rider paid.
\end{solution}

\begin{solution}{exo:m2:distressed-sovereign-and-bank-capital-credit:7}
Holding out is then worth 46.26 even if no one tenders, more than the 34.85 of
tendering: it pays at any participation below the clause's 75\%. A stronger legal
position widens the band to everything below the threshold, so only the clause stops
free-riding.
\end{solution}

\begin{solution}{exo:m2:distressed-sovereign-and-bank-capital-credit:8}
The ranking applies in a liquidation. An AT1 bond is also written down or converted by
its own terms and by the authority at the \omterm{def:m2:distressed-sovereign-and-bank-capital-credit:pon}{point of non-viability}, while the bank is
still operating, and nothing in those terms requires the shareholders to be wiped out
first; in Credit Suisse's case the AT1 bonds went to zero while shareholders received
CHF~3 billion of UBS shares. The EU authorities have said they would apply the
liquidation order; the Swiss case is before the courts.
\end{solution}

\begin{solution}{pb:m2:distressed-sovereign-and-bank-capital-credit:1}
\textbf{1.} 34.85 with the cash, 29.85 without.
\textbf{2.} 50\% of face; 65.2\% in present value.
\textbf{3.} 45.29 at 6\%, 27.76 at 12\%.
\textbf{4.} Because the new bonds pay a low coupon for fifteen years: their value, and
so the creditors' real loss, depends on the market's view of the sovereign after the
exchange.
\textbf{5.} 34.85 million: 29.85 million of new bonds and 5 million in cash.
\textbf{6.} Paid next year: $100/1.09 = 91.74$; the lawsuit: $130/1.09^6 = 77.51$; in
default: 15; the branch without early payment is worth $0.3 \times 77.51 + 0.7 \times
15 = 33.75$.
\textbf{7.} 33.75; 34.73; 38.97.
\textbf{8.} At 60.9\%.
\textbf{9.} The clause binds it to the exchange: it gets the new bonds, 29.85, without
the cash.
\textbf{10.} 58.72.
\textbf{11.} Not stably: at any participation in that band each creditor prefers to hold
out, which lowers participation; creditors who expect 75\% to be reached tender to keep
the cash, which raises it. Expectations decide which way it goes.
\textbf{12.} The cash is the difference between tendering and being bound, 5 per 100: it
is what a holdout loses by waiting once the clause binds.
\textbf{13.} By buying more than 25\% of the aggregated principal, which is far more than
a quarter of one series; aggregation makes a blocking stake much more expensive.
\textbf{14.} They were paid in full as their bonds matured, such as EUR~435 million in
May 2012, because the Greek law could not reach them and they were few.
\textbf{15.} In July 2014 they were not paid, because the courts barred payment on the
exchanged bonds unless the holdouts were paid in full at the same time.
\textbf{16.} It depends on the sovereign's cash, politics and legal exposure, and on how
many others hold out, all of which change during the restructuring.
\textbf{17.} A matter of judgement: holdouts enforce contracts and discipline debtors,
but they can block relief that most creditors accept; aggregated clauses let a large
majority decide.
\textbf{18.} Little is liquid: CDS on the sovereign, if any settled on the default,
are gone; the fund can hedge the exit yield with the sovereign's other traded bonds or
related currencies, but the legal outcome is unhedgeable.
\textbf{19.} Named result: \emph{the free-rider band}: holding out pays for participation
between 60.9\% and 75\%; at 75\% the clause binds the holdout to the offer, worth 29.85
against 34.85 for tendering.
\textbf{20.} They stop a minority from free-riding on a restructuring that a
qualified majority has accepted.
\end{solution}

\begin{solution}{iq:m2:distressed-sovereign-and-bank-capital-credit:1}
A company can file for bankruptcy, where a court imposes a plan on classes of
creditors and can liquidate assets; a sovereign cannot, has few seizable assets, and
restructures by \omterm{def:m2:distressed-sovereign-and-bank-capital-credit:distressed}{exchange offers} and contract terms. Holdouts can sue but struggle to
collect, and the outcome depends on governing law, \omterm{def:m2:distressed-sovereign-and-bank-capital-credit:cac}{collective action clauses} and
politics.
\iqlookfor{no court, contracts, holdouts.}
\end{solution}

\begin{solution}{iq:m2:distressed-sovereign-and-bank-capital-credit:2}
A perpetual bank bond counted as going-concern capital, with discretionary coupons,
written down or converted when a capital ratio trigger is hit or when the authority
declares the \omterm{def:m2:distressed-sovereign-and-bank-capital-credit:pon}{point of non-viability}. Risks: coupon cancellation, the bank not calling it
at the first call date, write-down or conversion, and the authorities' choice of loss
order, as 2023 showed.
\iqlookfor{coupons, extension, triggers, authority discretion.}
\end{solution}

\begin{solution}{iq:m2:distressed-sovereign-and-bank-capital-credit:3}
Value the new instruments at an exit yield estimated from comparable restructured
credits, add cash and any sweeteners, and compare with the recovery from holding out,
a probability tree over being paid, suing and staying in default, discounted
consistently; the choice depends on expected participation and the legal terms.
\iqlookfor{exit yield, holdout alternative, participation.}
\end{solution}

\begin{solution}{iq:m2:distressed-sovereign-and-bank-capital-credit:4}
Each creditor prefers others to accept the loss while it keeps its full claim, so
participation can stall. \omterm{def:m2:distressed-sovereign-and-bank-capital-credit:cac}{Collective action clauses} let a qualified majority bind the
minority; aggregated clauses count across all bonds, so a blocking stake in one bond
no longer suffices.
\iqlookfor{free-riding and the binding vote.}
\end{solution}

\begin{solution}{iq:m2:distressed-sovereign-and-bank-capital-credit:5}
That contractual terms and emergency law, not only the liquidation order, decide who
loses; that authorities can differ, as the EU's statement showed; and that the legal
basis can be challenged for years. Investors may now demand a premium for that
jurisdictional risk.
\iqlookfor{contract versus hierarchy, jurisdiction, legal risk.}
\end{solution}

\begin{solution}{iq:m2:distressed-sovereign-and-bank-capital-credit:6}
Parse each bond's prospectus into a record: governing law, pari passu wording, clause
type (series-by-series, two-limb, single-limb), thresholds, quorum and disenfranchised
holdings; keep outstanding amounts current from issuance and buyback data; compute
the votes needed under each procedure and the stake that would block each series;
flag changes when new bonds are issued.
\iqlookfor{structured contract data and vote arithmetic.}
\end{solution}
